\section{Related Work and Positioning}\label{sec:related}

\paragraph{Computational social choice.}
SAT and automated reasoning have long been used to discover and verify finite
social-choice impossibilities and to synthesize voting rules
\citep{tanglin2009computer,geist2011automated,brandtgeist2016strategyproof}.
The present contribution is not a new use of SAT for social choice and does not
claim a faster solver.  It studies a downstream explanatory question: after a
finite impossibility comparison is available, which aspects of a raw repair
presentation survive transfer to a neighboring case?

\paragraph{Finite-to-family transfer.}
Automated social-choice work has shown that finite instances can support wider
theorems under suitable transfer arguments, while integrity-constraint lifting
also illustrates that transfer conditions require care
\citep{geist2011automated,grandiendriss2013}.  M2.1 deliberately separates that
semantic question from the representation of repairs.  Even when impossibility
status is already legitimized by another layer, neither clause-multiset
equivalence nor repair-family equality follows automatically.

\paragraph{Diagnosis, MUS/MCS, and assumption systems.}
Minimal conflicts and minimal repairs are central to model-based diagnosis and
assumption-based truth maintenance \citep{dekleer1986atms,reiter1987}.
The Option D enumeration is adjacent to this tradition: lever removals are
tested for satisfiability and filtered by inclusion minimality.  We do not
propose a new MUS/MCS algorithm.  Our narrower point is that a minimal repair
family depends on the exposed lever vocabulary.  A bundled lever and its
fine-grained split can yield different raw singleton repairs even when their
clauses coincide.

\paragraph{Constraint hierarchies and optimization.}
Constraint hierarchies resolve inconsistency through preferences among
constraints \citep{borning1992constraint}; MaxSAT and related methods optimize
over soft clauses or weighted violations \citep{li2009maxsat}.  The transport
map $\pi$ in this study is neither a priority relation nor an objective
function.  It is an audit contract aligning two repair vocabularies.  The
question is representational equality of repair families, not the computation
of one preferred optimum.

\paragraph{Certified SAT and mechanized social choice.}
Proof assistants and certified clausal proof formats reduce trust in solver
internals \citep{moura2021lean4,cruzfilipe2017lrat,wetzler2014drat}.
Mechanized social-choice results provide a complementary semantic trust layer
\citep{nipkow2009social}.  The wider Sen atlas separates proof and artifact
layers.  This paper does not add a new Lean theorem or a new certificate.
Instead, it applies the same separation discipline to an explanatory layer:
status, clause structure, and repair presentation are recorded as distinct
claims with distinct evidence.

\paragraph{Encoding sensitivity.}
SAT encodings of one semantic problem can vary greatly in auxiliary variables,
propagation behavior, and proof structure.  Such variation is often judged by
solver performance.  Here it has a different role.  The Option C selector
variables are semantically motivated by witness choice, but their presence
breaks the exact clause-multiset control used by the base-case repair argument.
The relevant sensitivity is therefore not merely runtime sensitivity; it is
sensitivity of the admissible explanation claim to the encoding interface.
